\chapter{Plán realizace a ověření}
\label{chap:realizace}

\section{Nákup a aktualizovaný rozpočet}

Nákupní seznam jsem sjednotil s aktuálním výběrem komponent. Požadavek na
omezení počtu dodavatelů řeším dvěma skupinami: RPiShop pro výpočetní uzel,
IMU a rádia, ArduSimple pro oba přijímače a antény. ArduSimple dodává ze
Španělska; před objednávkou ověřím konečný daňový režim, dostupnost a dopravu.
\cite{ardusimple-faq}
Názvy položek v tabulkách odkazují přímo na veřejné produktové stránky.

\subsection{RPiShop: výpočetní část, IMU a LoRa}

\begin{table}[htbp]
  \centering
  \small
  \caption{Vybrané položky RPiShopu; ceny v Kč včetně DPH}
  \label{tab:rpishop}
  \begin{tabularx}{\textwidth}{@{}rYrr@{}}
    \toprule
    Ks & Položka & Kč/ks & Celkem \\
    \midrule
    1 & \href{https://rpishop.cz/535838/raspberry-pi-5-2gb-ram/}
      {Raspberry Pi 5, 2~GB (RPI502)} & 1\,629 & 1\,629 \\
    1 & \href{https://rpishop.cz/chladice-pro-raspberry-pi-5/6496-raspberry-pi-5-active-cooler.html}
      {Raspberry Pi 5 Active Cooler (OFI118)} & 149 & 149 \\
    1 & \href{https://rpishop.cz/502557/raspberry-pi-m-2-hat/}
      {Raspberry Pi M.2 HAT+ (OFI134)} & 349 & 349 \\
    4 & \href{https://rpishop.cz/595653/xiao-esp32s3-wio-sx1262-kit-pro-meshtastic-lora/}
      {XIAO ESP32S3 + Wio-SX1262 Kit (SEE319)} & 339 & 1\,356 \\
    1 & \href{https://rpishop.cz/652161/adafruit-9osy-imu-modul-icm-20948-se-stemma-qt-qwiic/}
      {Adafruit ICM-20948, 9osá IMU (ADA076)} & 449 & 449 \\
    \midrule
    & \textbf{Celkem s DPH} & & \textbf{3\,932} \\
    & Přepočet bez 21\,\% DPH & & 3\,249,59 \\
    \bottomrule
  \end{tabularx}
\end{table}

Ceny Raspberry Pi, chladiče a HAT+ vycházejí z podkladu z 24.~9.~2026.
LoRa a IMU jsou aktualizovány k 30.~9.~2026.
\cite{rpishop-pi5-2gb,rpishop-active-cooler,rpishop-m2hat,rpishop-xiao-wio,rpishop-icm20948}
Nová IMU nahrazuje původně uvažovanou položku od Farnellu a čtyři
XIAO/Wio sady nahrazují dvojici Waveshare. Staré i nové varianty se proto
do rozpočtu nesčítají.

Pro aktivní spoj jsou potřebné dvě rádiové sady. Dvě další jsou v počtu
čtyř zahrnuty pouze jednou a plánuji je pro vývoj a rezervu. Nepočítám se
čtyřmi sadami pro rover a dalšími čtyřmi pro základnu.

\subsection{ArduSimple: GNSS rover a přenosná základna}

\begin{table}[htbp]
  \centering
  \small
  \caption{Vybrané položky ArduSimple; ceny bez DPH}
  \label{tab:ardusimple}
  \begin{tabularx}{\textwidth}{@{}rYrr@{}}
    \toprule
    Ks & Položka & EUR celkem & Přepočet Kč \\
    \midrule
    1 & \href{https://www.ardusimple.com/product/simplertk3b-x5/}
      {simpleRTK3B Pro / mosaic-X5} & 569 & 13\,872,22 \\
    2 & \href{https://www.ardusimple.com/product/u-blox-gnss-tripleband-antenna-ann-mb2-00/}
      {u-blox ANN-MB2-00, rover a základna (61 EUR/ks)} & 122 & 2\,974,36 \\
    1 & \href{https://www.ardusimple.com/product/simplertk2b/}
      {simpleRTK2B Budget / ZED-F9P} & 172 & 4\,193,36 \\
    \midrule
    & \textbf{Celkem bez DPH} & \textbf{863} & \textbf{21\,039,94} \\
    \bottomrule
  \end{tabularx}
\end{table}

Podkladem jsou ceny přijímačů z nákupního seznamu z 24.~9.~2026 a cena
ANN-MB2-00 ověřená 30.~9.~2026.
\cite{ardusimple-x5,ardusimple-annmb2,ardusimple-f9p}
Pro srovnatelnost zachovávám kurz ČNB 24,380~Kč/EUR z 23.~9.~2026,
nikoli tvrzení o kurzu platném v den objednávky.\cite{cnb-rates}
Obě antény ANN-MB2-00 mají v ceně 5m kabel se SMA konektorem. Jeden kus
je určen pro rover a druhý pro základnu; nejde o dvouanténní rover.

\subsection{Mouser jako alternativa nákupní cesty}

Vedle přímé objednávky u ArduSimple jsem zvažoval přesun obou GNSS desek
a obou antén k distributorovi Mouser. Jde o alternativního prodejce
stejné koncepce, nikoli o druhou současnou objednávku. Předběžný rozdíl
jsem pro komunikaci se školitelem odhadl přibližně na 200~USD v neprospěch
Mouseru. Tento pracovní odhad není závaznou nabídkou a do tabulky
v korunách jej nepřičítám; případná změna vyžaduje novou kalkulaci všech
položek ve shodné měně a daňovém režimu.

Mouser může dávat smysl při jednodušším schvalování dodavatele univerzitou
nebo výhodnějším potvrzeném termínu. Samotná doména českého obchodu však
neprokazuje místo expedice. Před rozhodnutím bych ověřil shodu objednacích
kódů, hardwarovou revizi, příslušenství, sklad a dodací podmínky.
Současný nákupní podklad ponechávám rozdělený na \textbf{RPiShop a
ArduSimple}; přímá expedice ArduSimple ze Španělska odpovídá preferenci
evropského zdroje.\cite{ardusimple-faq}

\subsection{Součet a hranice kalkulace k 30. září 2026}

\begin{table}[htbp]
  \centering
  \small
  \caption{Souhrn oceněné elektroniky; nejsou započteny dosud nevybrané instalační díly}
  \label{tab:budget-2026}
  \begin{tabularx}{\textwidth}{@{}Yr@{}}
    \toprule
    Skupina & Kč bez DPH \\
    \midrule
    RPiShop: 3\,932 / 1,21 & 3\,249,59 \\
    ArduSimple: 863 EUR při 24,380 Kč/EUR & 21\,039,94 \\
    \textbf{Vybraná elektronika celkem} & \textbf{24\,289,53} \\
    Cílová částka & 25\,000,00 \\
    \textbf{Rezerva před dopravou a instalací} & \textbf{710,47} \\
    \midrule
    Orientační celek s 21\,\% DPH & 29\,390,33 \\
    \bottomrule
  \end{tabularx}
\end{table}

Výpočet vychází z nezaokrouhleného přepočtu částky RPiShopu; součet samostatně
zaokrouhlených řádků se může lišit o haléř. Pro stručnou komunikaci používám
\textbf{24~290~Kč bez DPH} a rezervu přibližně \textbf{710~Kč}.
Kontrolní cena s DPH je orientační sjednocený přepočet, ne závazná nabídka
ani určení daňového režimu faktury univerzity.

SSD dodám vlastní. Pro první laboratorní etapu počítám s týmovým
notebookem základny; předpoklad jeho dostupnosti musím potvrdit.
U SSD ověřím NVMe provedení, formát 2230/2242 a kompatibilitu s HAT+.
Notebook mohu při cílovém samostatném provozu vypustit až po doplnění
a ověření lokálního raw loggeru. Košík tedy není kompletní cenou
autonomní základny včetně veškerého příslušenství.

Částka \textbf{nezahrnuje} dopravu, DC/DC zdroj, ochrany a energetický
doběh, USB a signálovou kabeláž, konektory, kryty, držáky, stativ ani
případné kvalitnější LoRa antény. Nezahrnuje ani samostatné úložiště
raw dat základny a jeho rozhraní. Tyto položky zatím nemají konečnou
specifikaci a cenu, proto je neoceňuji nulou. Celková cena funkční instalace
bude vyšší. Původní orientační kalkulace s jinými deskami a modemy už není
aktuálním podkladem objednávky.

\section{Etapy sestavení}

\subsection{Stolní prototyp}

Nejprve ověřím obě GNSS desky s anténami, Raspberry Pi s chlazením a SSD
a přímé SPI/INT připojení IMU. Laboratorní napájení musí mít dostatečnou
proudovou rezervu pro celou sestavu. V této etapě nechci návrh uzavírat do
finální krabičky, která by ztížila přístup k měřicím bodům.

Sběr proběhne nejprve bez online filtru. Po ověření souborových formátů a
časování přidám fúzi a vyzkouším současnou zátěž CPU, USB a NVMe. LoRa
vyvíjím nejprve s generovanými daty, potom se zaznamenaným RTCM a nakonec
s živou základnou. Předinstalovaný Meshtastic nahradím firmwarem, který
umožní kontrolovat rámce, čas ve vzduchu a stáří přenášených korekcí.

Následně přidám konfiguraci ZED-F9P přes UART a lokální webovou správu
ESP32. Odděleně ověřím souběh Wi-Fi, rádia a zápisu raw dat. Teprve když
základna uchová kompletní observace i po výpadku rádia a správně uzavře
záznam při vypnutí, přejdu k provozu bez notebooku.

\subsection{Doplňky nutné pro provoz}

\small
\begin{longtable}{@{}L{0.27\textwidth}L{0.33\textwidth}L{0.32\textwidth}@{}}
  \caption{Instalační položky mimo původní rozpočet a jejich ověření}\label{tab:bom-b}\\
  \toprule
  Položka & Účel & Co zbývá rozhodnout \\
  \midrule
  \endfirsthead
  \toprule
  Položka & Účel & Co zbývá rozhodnout \\
  \midrule
  \endhead
  Napájení roveru & DC/DC, jištěný výstup PMU, filtrace a ochrany & nastavení PMU, příkon, přepětí a bezpečné vypnutí \\
  USB kabely & GNSS a XIAO k hostitelům & konektory, datové provedení, délka a jištění; kit XIAO kabel nemá \\
  SPI, INT a PPS kabeláž & signály a společná zem & dostupné GPIO pod HAT+, délky a úrovně PPS \\
  Kryt a držáky & ochrana a opakovatelné ustavení & vibrace, odvod tepla, magnetické rušení a montáž antén \\
  Základnový stativ & stabilní bod s volným horizontem & výška, zaměření a zajištění proti posunu \\
  Napájení základny & přijímač, rádio a hostitel & výdrž při celé disciplíně a záznam poklesu napětí \\
  Lokální úložiště základny & raw log pro PPK bez notebooku & například microSD, rozhraní a piny nekolidující s Wio; kapacita, výdrž zápisu a zotavení \\
  Vodivá podložka antény & ověření montážní konfigurace & materiál, rozměry, mechanické uchycení a A/B test krytu \\
  LoRa antény pro provoz & rádiová rezerva a vhodná montáž & nejprve test dodaných FPC; případné náhrady včetně zisku a konektorů \\
  Převodník pro FSYNC, volitelný & případné rozšíření synchronizace & \SI{1.8}{V} úroveň a ověření významu synchronizačního záznamu \\
  \bottomrule
\end{longtable}
\normalsize

Základna bude mít ZED-F9P, anténu ANN-MB2-00 a jednu aktivní
rádiovou sadu. Raw data nejprve zaznamená existující notebook, cílově
doplněné místní úložiště. Rover bude mít X5 s ANN-MB2-00,
IMU, Raspberry Pi a druhou aktivní rádiovou sadu. Tento rozpis určuje použití
již oceněných komponent, nikoli další nákup stejných kusů.

\section{Kontroly před objednávkou}

U ArduSimple chci potvrdit konkrétní variantu obou přijímačů, dostupnost
a obsah dodávky. Pro X5 je zásadní současná frekvence raw observací a
skutečného RTK řešení při vybrané konfiguraci, potřebná licence a firmware.
Dále ověřím vyvedení PPS, jeho elektrické parametry a vhodnost antény.
U základnové desky potvrdím revizi, fyzické rozměry a nastavení IOREF
pro UART. Držák neuzavřu podle obrázku jiné revize. U antén ověřím pevný
5m kabel a místo pro jeho bezpečné vedení ve vlastním krytu.

U RPiShopu ověřím dostupnost jedné desky ADA076 a všech čtyř sad SEE319.
U ICM-20948 beru v úvahu EOL a neplánuji dostupnost náhrad podle aktuálního
jediného skladového kusu. Zkontroluji zapojení SPI, INT a odlišné napětí
FSYNC podle dokumentace Adafruit. Konektor STEMMA QT není náhradou za SPI
kabeláž.

Před potvrzením nákupu je potřeba doplnit dopravu a neoceněné povinné díly
a porovnat celkovou částku s rozpočtem. Volba XIAO/Wio znamená také vlastní
vývoj firmwaru a ověření rádiového režimu; nejde o nákup dvou již hotových
a odzkoušených RTCM modemů.

\section{Implementační plán}

\begin{table}[htbp]
  \centering
  \small
  \caption{Pracovní balíčky a jejich ověřitelné výstupy}
  \label{tab:plan}
  \begin{tabularx}{\textwidth}{@{}L{0.10\textwidth}L{0.29\textwidth}Y@{}}
    \toprule
    Balíček & Obsah & Výstup \\
    \midrule
    WP1 & Rozhraní a nákup & úplný soupis včetně napájení, dopravy a potvrzení dostupnosti \\
    WP2 & IMU a časování & SPI/FIFO/INT ovladač, PPS, odhad epoch a protokol nejistoty času \\
    WP3 & Záznam na Pi OS & tři proudy, metadata, diagnostika ztrát a bezpečné ukončení \\
    WP4 & Základna a LoRa & raw log, UART konfigurace, webová správa, RTCM most a ověřený rádiový režim \\
    WP5 & Online ESKF & poloha, rychlost, orientace, biasy, kovariance a odmítání rušení \\
    WP6 & Postprocessing & převod observací, PPK a reprodukovatelné vyhlazení \\
    WP7 & Kalibrace a montáž & lever arm, boresight, teplota, vibrace a magnetická kalibrace \\
    WP8 & Validace & výsledky zkoušek proti kritériím tabulky~\ref{tab:acceptance} \\
    \bottomrule
  \end{tabularx}
\end{table}

Online fúzi a definitivní kryt zařadím až po ověření raw loggeru. Změna
senzoru nebo časového modelu tak neznehodnotí již pořízené záznamy.

\section{Ověřovací a kalibrační plán}

\subsection{Časování a úplnost záznamu}

Na osciloskopu nebo logickém analyzátoru budu současně sledovat PPS, INT
a SPI přenosy. Porovnám fyzické hrany s kernelovými časovými značkami,
periodou měření a obsahem FIFO. Test má trvat nejméně 30 minut v klidu
i se zatíženým CPU, USB, rádiem a diskem. Vyhodnotím zpoždění, jitter,
přetečení, návaznost počtu vzorků a reakci na odpojení senzoru.

Zvlášť ověřím akcelerometr, gyroskop a magnetometr. Nový magnetický vzorek
musí být odlišen od staré hodnoty a musí mít vlastní odhad času a stav
platnosti. Ztráta jedné události nesmí vést k nepozorovanému posunu všech
následujících epoch. Pokud nedokážu přiřazení obnovit, logger uzavře
časově souvislý úsek a označí mezeru.

\subsection{Inerciální a magnetická kalibrace}

Provedu šestipolohovou kontrolu akcelerometru, kontrolovaný rotační pohyb
a statický záznam pro posouzení šumu a Allanova rozptylu. Porovnám bias po
opakovaném zapnutí a při různých teplotách. Tyto zkoušky nejsou automaticky
úplnou teplotní kalibrací, ale mohou odhalit nevhodné nastavení nebo kus.

Magnetometr nejprve zkalibruji mimo vozidlo a poté ověřím v konečné montáži.
Test zahrne více orientací, vypnutou a zapnutou výkonovou elektroniku
a různé proudové zatížení. Vyhodnotím velikost i směr pole a rezidua po
hard-iron/soft-iron korekci. V datech uchovám původní i kalibrované hodnoty
a verzi kalibrace. Pokud rušení za jízdy nelze spolehlivě rozpoznat,
magnetické aktualizace do filtru nezařadím.
Součástí testu bude také vliv magnetu GNSS antény, šroubů a držáků;
stálou odchylku i změnu pole při provozu vyhodnotím odděleně.

\subsection{GNSS a rádiový spoj}

Rover a základnu postavím na stabilní body s otevřeným horizontem. Změřím
skutečnou četnost observací a PVT, dobu do FIX, stáří korekcí a datový objem
RTCM. Stejný proud porovnám přes kabel, LoRa a případně NTRIP. Tím oddělím
chybu rádiové vrstvy od problému samotného GNSS.

Pro přesnost použiji lokální referenci navázanou na základnu: například
nezávisle zaměřené body a známé vzdálenosti. Samostatně uvedu vodorovné
a svislé chyby, jejich rozdělení a nejistotu reference. Prosté odčítání
průměrné polohy měří opakovatelnost, nikoli samo o sobě přesnost.
Při opakování zkoušky zachovám počátek ENU a montáž; po změně základny
zdokumentuji převod do společné soustavy.

Kryt a podložku antény porovnám při zachování jejího referenčního bodu
a pokud možno obdobného výhledu i družicové geometrie. Test zopakuji
s aktivním Wi-Fi, LoRa a DC/DC. Zaznamenám C/N0, FIX, rozptyl i odchylku
od reference. Jediný krátký test nepovažuji za obecný důkaz lepší antény.

U LoRa zaznamenám konfiguraci, počty zpráv a fragmentů, chyby integrity,
ztráty, RSSI/SNR, latenci a čas ve vzduchu. Provedu test až do požadované
vzdálenosti přibližně \SI{1}{km} a projedu trať s reálnou výškou antény
a orientací vozidla. Nevyžaduji fyzikálně nemožnou záruku nulové rádiové
ztrátovosti; požaduji její měřitelnost a přijatelné stáří použitých korekcí.

Řízeně přeruším spoj a ověřím nepřerušený raw log na obou koncích, označení
zhoršené kvality a obnovení RTK po návratu korekcí. Následně ověřím, že lze
jízdu zpracovat PPK i přes výpadek rádiového přenosu.

\subsection{Dynamika a napájení na vozidle}

Po zaměření montáže začnu pomalou jízdou, osmičkou, brzděním a slalomem.
Teprve potom přejdu k plné dynamice. U každé osy vyhodnotím maxima,
99,9.~percentil a četnost překročení 80\,\% rozsahu. Z naměřeného spektra
posoudím rezonance držáku, filtraci a potřebu jiného rozsahu.

Změřím teplotu Raspberry Pi a SSD, napájecí poklesy a vliv rádia i DC/DC
měniče na GNSS a magnetometr. Test řízeného vypnutí musí ověřit čitelnost
již uzavřených bloků a označení neúplné relace. Energetický doběh se musí
navrhnout podle skutečné doby potřebné k bezpečnému ukončení.

\small
\begin{longtable}{@{}L{0.23\textwidth}L{0.38\textwidth}L{0.31\textwidth}@{}}
  \caption{Navržená akceptační kritéria; nejde o dosažené výsledky}\label{tab:acceptance}\\
  \toprule
  Zkouška & Kritérium / pracovní cíl & Doklad \\
  \midrule
  \endfirsthead
  \toprule
  Zkouška & Kritérium / pracovní cíl & Doklad \\
  \midrule
  \endhead
  IMU, 30 minut & žádné nevysvětlené mezery ani přetečení při plánované zátěži & počty dat, událostí a příznaky FIFO \\
  Vazba epoch & změřená nejistota; pracovní cíl pod \SI{100}{\micro\second} & externí záchyt hran, rezidua a zpoždění filtrů \\
  Magnetometr & rozlišené nové vzorky a ověřené odmítnutí rušených dat & kalibrační protokol a srovnání fúze s/bez kompasu \\
  GNSS frekvence & dosažena dodavatelem potvrzená raw i RTK četnost & rozdíly časů SBF epoch \\
  Statické RTK & pracovní cíl 95\,\% epoch FIX po inicializaci na otevřeném bodě & stav, korekce a referenční souřadnice \\
  Místní přesnost & cíle 2 cm vodorovně a 3 cm svisle posouzené proti nezávislé referenci & chyby, jejich statistika, nejistota reference a společný rámec ENU \\
  Samostatná základna & úplný raw log a přenos při souběhu s webovou správou & srovnání s USB záznamem, zátěž a zkouška restartu \\
  LoRa na trati & změřená ztrátovost, omezené stáří korekcí a vyhovující rádiový režim & čas ve vzduchu, latence, RSSI/SNR a RTK stav \\
  Saturace & žádná při běžné jízdě; rezerva k 80\,\% rozsahu & maxima, percentily a označené špičky \\
  Výpadek spojení & raw log pokračuje a korekce se po návratu obnoví & události a návaznost relací \\
  Výpadek napájení & uzavřené bloky čitelné, neúplnost zjistitelná & kontrolní součty a opakovaný start \\
  Trajektorie & rozdíly online/PPK/smoothing mají vysvětlitelné příčiny & rezidua, reference a rozdělení podle kvality \\
  \bottomrule
\end{longtable}
\normalsize

Kritérium 95\,\% FIX se týká otevřeného referenčního testu, nikoli garance
pro celou závodní trať. Pokud časování nedosáhne cíle, musím skutečnou
nejistotu promítnout do fúze a rozhodnout o změně synchronizace.

\section{Rizika a rozhodovací milníky}

\small
\begin{longtable}{@{}L{0.32\textwidth}L{0.60\textwidth}@{}}
  \caption{Hlavní rizika aktualizovaného návrhu}\label{tab:risks}\\
  \toprule
  Riziko & Opatření \\
  \midrule
  \endfirsthead
  \toprule
  Riziko & Opatření \\
  \midrule
  \endhead
  Jiná RTK četnost X5 než očekávaná & potvrzení varianty a konfigurace; kontrola skutečných epoch před montáží \\
  Drift nebo vyšší šum ICM-20948 & statické testy, kalibrace a porovnání výpadků GNSS; případný návrat k jiné IMU \\
  Magnetické rušení & měření ve vozidle, kalibrace, detekce odlehlých dat a možnost kompas nepoužívat \\
  EOL a nedostupnost další desky & potvrdit aktuální kus; neodvozovat dlouhodobou dostupnost z jedné nabídky \\
  Jitter Linuxu nebo ztráta FIFO & přerušení, FIFO, zátěžová zkouška, časová nejistota a značení přerušených úseků \\
  RTCM se nevejde do LoRa režimu & změřit objem a airtime; redukovat zprávy jen při zachování RTK; záložní přenos \\
  Firmware rádia není hotový modem & samostatné testy rámcování, fragmentace, resetu, stáří dat a rádiové shody \\
  Anténa nebo držák zhorší měření & montážní A/B test, vhodná zemní plocha, zaměření a kontrola vibrací \\
  Nový survey-in posune porovnávané jízdy & uložit souřadnice a identifikátor ustavení; společný počátek nebo doložená transformace \\
  Samostatná základna ztratí raw data & oddělit logger od rádiové fronty; ověřit úložiště před odpojením notebooku \\
  Rozpočet není úplný & doplnit napájení, instalační díly a dopravu před schválením celkové objednávky \\
  \bottomrule
\end{longtable}
\normalsize

Prvním milníkem je úplný raw logger se změřenou časovou vazbou. Druhým je
současně funkční základna, RTK rover, ověřený přenos korekcí a reprodukovatelné
PPK. Třetím je vozidlová validace včetně magnetického rušení, saturace,
teploty a napájení. Teprve po těchto krocích má smysl uzavírat návrh finálního
krytu a případné vlastní desky.

\section{Dílčí závěr výběru GNSS/IMU}

Pro první realizaci jsem vybral \textbf{simpleRTK3B Pro s mosaic-X5},
jedinou přenosnou základnu \textbf{simpleRTK2B Budget se ZED-F9P},
dvě antény \textbf{ANN-MB2-00} (po jedné pro rover a základnu),
\textbf{Adafruit ICM-20948} a \textbf{Raspberry Pi 5 2~GB} s chlazením,
M.2 HAT+ a vlastním SSD. Přímý korekční spoj mají zajistit dvě ze čtyř
objednávaných sad \textbf{XIAO ESP32S3 + Wio-SX1262}.

Původní šestiosou IMU jsem změnil s ohledem na možnost využít kompas,
dostupnost a sdružení nákupu. Přínos této změny není automatický:
magnetometr musí projít kalibrací a zkouškou rušení, zatímco šum a drift
inerciálních kanálů musí odpovídat záměru práce. EOL čipu omezuje
perspektivu opakované výroby. Pro aktuální prototyp tuto nevýhodu přijímám,
ale v dokumentaci ji nezamlčuji.

Raspberry Pi bude číst IMU přímo přes SPI a INT, bez samostatného
časovacího STM32. GNSS PPS a GPIO události vytvoří základ společné
časové osy, jejíž nejistotu ověřím měřením. ESP32 zůstává součástí rádiových
sad a u základny má navíc zajistit konfiguraci a místní webovou správu.
Pro samostatný provoz musím doplnit úložiště raw dat. Vlastní firmware, legální rádiová konfigurace
a změřená propustnost jsou nezbytné realizační úkoly.

Aktuální oceněná elektronika vychází na \textbf{24~289,53~Kč bez DPH},
tedy přibližně \textbf{24~290~Kč}. Rezerva do cílové částky je asi
710~Kč před dopravou a dosud neoceněnou instalací. Nákup je soustředěn
u dvou dodavatelů. Úplný rozpočet bude možné uzavřít až po volbě napájení,
kabeláže a mechanických dílů.

Dvě ANN-MB2 volím pro společné montážní řešení a nižší cenu; jejich
přesnost v krytu teprve ověřím. Prioritou je místní trajektorie vůči
základně, nikoli centimetrové umístění na mapě. Volba vývojové desky
místo Holybro vychází z přístupu k rozhraním a vlastní konstrukce,
ne z předpokladu přesnějšího GNSS čipu. Mouser ponechávám jako
alternativu nákupu, nikoli dalšího současného dodavatele.

Zachování tří měřicích proudů umožní porovnávat online filtr s pozdějším
PPK a vyhlazením. Za nejdůležitější výsledek přípravné fáze považuji
ověřitelný návrh a jasné experimentální otázky; skutečnou přesnost
trajektorie bude možné uvést až po realizaci a validaci.

Tento výběr představuje navigační část širšího měřicího systému.
Kapitola~\ref{chap:integrace} rozšiřuje návrh o senzoriku vozidla,
centrální ADC, dvě rozhraní CAN pro auto, napájení a montáž. Nové komponenty nejsou
automaticky zahrnuty do zde uvedeného cenového snímku GNSS/IMU sestavy.
